\end{proof}

\begin{lemma}
\label{lemma-conormal-sequence-H1-regular-ideal}
Let $A$ be a ring. Let $I \subset J \subset A$ be ideals.
Assume that $J/I \subset A/I$ is a $H_1$-regular ideal.
Then $I \cap J^2 = IJ$.
\end{lemma}

\begin{proof}
Follows immediately from Lemma \ref{lemma-conormal-sequence-H1-regular}
by localizing.
\end{proof}







\section{Local complete intersection maps}
\label{section-lci}

\noindent
We can use the material above to define a local complete intersection
map between rings using presentations by (finite) polynomial algebras.

\begin{lemma}
\label{lemma-koszul-independence-presentation}
Let $A \to B$ be a finite type ring map. If for some presentation
$\alpha : A[x_1, \ldots, x_n] \to B$ the kernel $I$ is a Koszul-regular ideal
then for any presentation $\beta : A[y_1, \ldots, y_m] \to B$ the kernel
$J$ is a Koszul-regular ideal.
\end{lemma}

\begin{proof}
Choose $f_j \in A[x_1, \ldots, x_n]$ with $\alpha(f_j) = \beta(y_j)$
and $g_i \in A[y_1, \ldots, y_m]$ with $\beta(g_i) = \alpha(x_i)$.
Then we get a commutative diagram
$$
\xymatrix{
A[x_1, \ldots, x_n, y_1, \ldots, y_m]
\ar[d]^{x_i \mapsto g_i} \ar[rr]_-{y_j \mapsto f_j} & &
A[x_1, \ldots, x_n] \ar[d] \\
A[y_1, \ldots, y_m] \ar[rr] & & B
}
$$
Note that the kernel $K$ of $A[x_i, y_j] \to B$ is equal to
$K = (I, y_j - f_j) = (J, x_i - g_i)$. In particular, as
$I$ is finitely generated by
Lemma \ref{lemma-quasi-regular-ideal-finite}
we see that $J = K/(x_i - g_i)$ is finitely generated too.

\medskip\noindent
Pick a prime $\mathfrak q \subset B$. Since
$I/I^2 \oplus B^{\oplus m} = J/J^2 \oplus B^{\oplus n}$
(Algebra, Lemma \ref{algebra-lemma-conormal-module})
we see that
$$
\dim J/J^2 \otimes_B \kappa(\mathfrak q) + n =
\dim I/I^2 \otimes_B \kappa(\mathfrak q) + m.
$$
Pick $p_1, \ldots, p_t \in I$ which map to a basis of
$I/I^2 \otimes \kappa(\mathfrak q) = I \otimes_{A[x_i]} \kappa(\mathfrak q)$.
Pick $q_1, \ldots, q_s \in J$ which map to a basis of
$J/J^2 \otimes \kappa(\mathfrak q) = J \otimes_{A[y_j]} \kappa(\mathfrak q)$.
So $s + n = t + m$. By Nakayama's lemma there exist $h \in A[x_i]$ and
$h' \in A[y_j]$ both mapping to a nonzero element of $\kappa(\mathfrak q)$
such that $I_h = (p_1, \ldots, p_t)$ in $A[x_i, 1/h]$ and
$J_{h'} = (q_1, \ldots, q_s)$ in $A[y_j, 1/h']$.
As $I$ is Koszul-regular we may also assume that $I_h$ is generated
by a Koszul regular sequence. This sequence must necessarily have length
$t = \dim I/I^2 \otimes_B \kappa(\mathfrak q)$, hence we see that
$p_1, \ldots, p_t$ is a Koszul-regular sequence by
Lemma \ref{lemma-independence-of-generators}.
As also $y_1 - f_1, \ldots, y_m - f_m$ is a regular sequence we
conclude
$$
y_1 - f_1, \ldots, y_m - f_m, p_1, \ldots, p_t
$$
is a Koszul-regular sequence in $A[x_i, y_j, 1/h]$
(see Lemma \ref{lemma-join-koszul-regular-sequences}).
This sequence generates the ideal $K_h$. Hence the
ideal $K_{hh'}$ is generated by a Koszul-regular sequence
of length $m + t = n + s$. But it is also generated by the sequence
$$
x_1 - g_1, \ldots, x_n - g_n, q_1, \ldots, q_s
$$
of the same length which is thus a Koszul-regular sequence by
Lemma \ref{lemma-independence-of-generators}.
Finally, by
Lemma \ref{lemma-truncate-koszul-regular}
we conclude that the images of $q_1, \ldots, q_s$ in
$$
A[x_i, y_j, 1/hh']/(x_1 - g_1, \ldots, x_n - g_n)
\cong A[y_j, 1/h'']
$$
form a Koszul-regular sequence generating $J_{h''}$. Since $h''$
is the image of $hh'$ it doesn't map to zero in $\kappa(\mathfrak q)$
and we win.
\end{proof}

\noindent
This lemma allows us to make the following definition.
